\section{A multiscale selector estimate}
\label{sec:selector}

The key estimate finds a single member of the family whose mass is
captured at every scale.  The family in this section need not be
increasing.  Recall that a family $\F\subseteq 2^X$ is
\emph{$p$-small} if there is a family $\G\subseteq 2^X$ such that every
$H\in\F$ contains some $G\in\G$ and
$\sum_{G\in\G}p^{|G|}\le 1/2$.

The argument combines the minimum-fragment method of Park and Pham
\cite[Definitions~2.3 and~2.5, Proposition~2.6]{ParkPham2024selector}
and its multiscale development by Pham
\cite[arXiv version~1, Section~2]{Pham2024} with the maximal-cutoff
argument of Bednorz, Martynek, and Meller~\cite[Lemma~3.6]{BMM2022}.
All the estimates needed here are proved below.

\begin{theorem}[Multiscale selector estimate]
\label{thm:multiscale-selector}
Let $X$ be a finite set, let $\F$ be a nonempty family of nonempty
subsets of $X$, and suppose that $\F$ is not $p$-small, where
$0<p\le 1$.  For each $H\in\F$, let $\lambda_H$ be a probability mass
function on $X$ supported on $H$.  Set $D=256$, and let $s\ge 0$ be an
integer such that
\[
  p\sum_{i=1}^s D^i\le \frac12.
\]
Independently color the elements of $X$ by a map
$a:X\to\{1,\ldots,s+1\}$, with color probabilities
\begin{equation}
\label{sel:color-probabilities}
  \pi_i=D^ip\quad(1\le i\le s),
  \qquad
  \pi_{s+1}=1-p\sum_{i=1}^sD^i.
\end{equation}
For a coloring $v$, write $A_i(v)=\{x\in X:v(x)\le i\}$.
With probability at least $9/10$, there is an $H\in\F$ such that
\begin{equation}
\label{sel:simultaneous-capture}
  \lambda_H(A_i(a))\ge 1-2^{-i}
  \qquad(1\le i\le s).
\end{equation}
In particular, with probability at least $9/10$, there is an
$H\in\F$ for which
\begin{equation}
\label{sel:average-color}
  \sum_{x\in X}a(x)\lambda_H(x)\le 2.
\end{equation}
\end{theorem}

For colorings where simultaneous capture fails, the proof will construct
covers from sets of vertices moved to earlier colors. We will sum these
costs against the probabilities of the failing colorings. The estimate
comes from locating the changed vertices within small sets determined
by the final coloring and the numbers moved between colors. The
following truncation observation supplies the size bound for these sets.

\begin{lemma}[Maximal truncation]
\label{sel:truncation}
Let $\lambda$ be a probability mass function on a finite set $X$,
let $A\subseteq X$, and let $m\ge 0$ and $0<d<1$.
Call $u\in[0,1]$ admissible if
\begin{equation}
\label{sel:admissible-cutoff}
  \{x:\lambda(x)>u\}\subseteq A,
  \qquad
  \sum_{x\in X}\min\{\lambda(x),u\}
       \bigl(\mathbf 1_A(x)-(1-d)\bigr)\ge mu.
\end{equation}
If there is an admissible cutoff, there is a largest one,
$\varepsilon$, and it satisfies
\[
  |\{x:\lambda(x)>\varepsilon\}|\le \frac{m}{d}.
\]
\end{lemma}

\begin{proof}
The inclusion in \eqref{sel:admissible-cutoff} is equivalent to
\[
  u\ge \max\bigl(\{\lambda(x):x\notin A\}\cup\{0\}\bigr).
\]
The remaining inequality is closed in $u$, since its two sides are
continuous.  Thus the nonempty set of admissible cutoffs is compact
and has a largest element $\varepsilon$.

Put $R=\{x:\lambda(x)>\varepsilon\}$.  If $\varepsilon=1$, then
$R=\varnothing$ and the assertion holds.  Otherwise, consider
\[
  f(u)=\sum_{x\in X}\min\{\lambda(x),u\}
          \bigl(\mathbf 1_A(x)-(1-d)\bigr)-mu.
\]
On a sufficiently short interval to the right of $\varepsilon$, the
only unsaturated terms are those indexed by $R$.  Each of these
points belongs to $A$, so the slope on that interval is
$d|R|-m$.  This statement also holds when $\varepsilon=0$ or when
$\varepsilon$ equals one of the weights: weights equal to the cutoff
remain constant when the cutoff increases.  An increase of the
cutoff preserves the inclusion in \eqref{sel:admissible-cutoff}.
If $d|R|-m$ were positive, it would also preserve $f(u)\ge 0$ for a
small increase, contradicting maximality.  Hence $d|R|\le m$.
\end{proof}

\begin{proof}[Proof of Theorem~\ref{thm:multiscale-selector}]
If $s=0$, the only color is $1$, and both conclusions are immediate.
Assume henceforth that $s\ge 1$, and put $d_i=2^{-i}$ for
$1\le i\le s$.  Call a coloring $a$ \emph{bad} if no $H\in\F$
satisfies \eqref{sel:simultaneous-capture}.  We will prove
$\P(a\text{ is bad})\le 1/10$.

\medskip
\noindent\textit{Feasible moves and a cover.}
We may move vertices to earlier colors, obtaining a coloring
$z\le a$ pointwise.  Write
\[
  U(a,z)=\{x:z(x)<a(x)\},
  \qquad
  m_i(a,z)=|A_i(z)\setminus A_i(a)|.
\]
Call such a move \emph{feasible} if there are $H\in\F$ and
$\varepsilon_1,\ldots,\varepsilon_s\in[0,1]$ such that, for every $i$,
\begin{equation}
\label{sel:feasibility}
\begin{split}
  R_i:=\{x:\lambda_H(x)>\varepsilon_i\}&\subseteq A_i(z),\\
  \sum_{x\in X}\min\{\lambda_H(x),\varepsilon_i\}
       \bigl(\mathbf 1_{A_i(z)}(x)-(1-d_i)\bigr)
       &\ge m_i(a,z)\varepsilon_i.
\end{split}
\end{equation}
The right side charges $\varepsilon_i$ for each vertex in
$A_i(z)\setminus A_i(a)$. Removing one of these vertices lowers the
truncated sum by at most $\varepsilon_i$ while reducing that charge by exactly $\varepsilon_i$;
hence the inequality survives. If the vertex has mass at most
$\varepsilon_i$, it lies outside $R_i$, so the inclusion survives as
well. We will apply this balance when delaying a vertex from color $i$
to color $i+1$.

If the unchanged move $z=a$ is feasible, then $a$ is not bad.
Indeed, in this case $m_i(a,a)=0$.  Replacing each truncated weight
by its original value adds only nonnegative terms to the left side
of \eqref{sel:feasibility}: any excess weight lies in
$R_i\subseteq A_i(a)$ and has coefficient $d_i$.
The resulting inequality is
$\lambda_H(A_i(a))-(1-d_i)\ge0$, for the same $H$ at every $i$.

Fix orders on the finite sets of colorings and members of $\F$.
For every $a$ and $H_0\in\F$, choose a feasible coloring $z$ with
$U(a,z)\subseteq H_0$ minimizing
\[
  L(a,z)=\sum_{x\in X}(a(x)-z(x)),
\]
breaking ties by the fixed order.  This minimum exists: the move
that sets $z(x)=1$ on $H_0$ and leaves other colors unchanged is
feasible with $H=H_0$ and all cutoffs zero.  Notice that the witness
$H$ in the minimization is otherwise unrestricted; it need not
equal $H_0$.

Let $\mathcal Z(a)$ be the set of chosen colorings, with repetitions
removed.  The changed sets $U(a,z)$, $z\in\mathcal Z(a)$, cover
$\F$, because each $H_0$ contains its chosen changed set.  Failure of
$p$-smallness therefore implies
\[
  \sum_{z\in\mathcal Z(a)}p^{|U(a,z)|}>\frac12.
\]
This remains true if distinct colorings have the same changed set:
the displayed sum is then at least the cost of the cover with its
repeated sets removed.  Define
$P(v)=\prod_{x\in X}\pi_{v(x)}$ for every coloring $v$.  It follows
that
\begin{equation}
\label{sel:cover-cost}
  \frac12\P(a\text{ is bad})
  \le
  \sum_{a\text{ bad}}P(a)
       \sum_{z\in\mathcal Z(a)}p^{|U(a,z)|}.
\end{equation}
Every changed set appearing on the right is nonempty, since a bad
coloring has no feasible unchanged move.

\medskip
\noindent\textit{A witness determined by the final coloring and the counts.}
Group the pairs $(a,z)$ on the right of \eqref{sel:cover-cost} by
$z$ and the array of counts
\[
  n_{ih}=|\{x:z(x)=i,\ a(x)=h\}|
  \qquad(1\le i<h\le s+1).
\]
We call this array the \emph{profile}.  Put
\[
  t_i=\sum_{h>i}n_{ih},
  \qquad t=\sum_{i=1}^s t_i=|U(a,z)|.
\]
The profile also determines every quantity in the feasibility tests
that depends on $a$, because
\begin{equation}
\label{sel:profile-crossings}
  m_i(a,z)=\sum_{j\le i<h}n_{jh}.
\end{equation}

Fix $z$ and a profile occurring in the sum.  Choose the first
$H\in\F$ that can witness \eqref{sel:feasibility} with these data.
Such a witness exists, and this choice uses only $z$ and the
profile.  For this fixed $H$, choose at every $i$ the largest
admissible cutoff in \eqref{sel:feasibility}.  The choices for the
different indices impose no restrictions on one another.
Lemma~\ref{sel:truncation} gives
\begin{equation}
\label{sel:large-weight-bound}
  |R_i|\le \frac{m_i(a,z)}{d_i}\le 2^i t.
\end{equation}
The resulting $H$, cutoffs, and sets $R_i$ are fixed for the entire
group.  In particular, they witness feasibility for every original
coloring $a$ in that group, irrespective of the member $H_0$ that
led to the choice of $z$.

For each such original coloring we claim that
\begin{equation}
\label{sel:changed-vertices}
  \{x:z(x)=i<a(x)\}\subseteq R_i
  \qquad(1\le i\le s).
\end{equation}
Suppose instead that a vertex $x$ on the left lies outside $R_i$.
Delay its color in $z$ from $i$ to $i+1$, obtaining $z'$.
Since $a(x)>i$, we still have $z'\le a$.  Choose an $H_0$ for which
$z$ was a minimum.  The changed set of $z'$ is contained in that of
$z$, so it is still contained in $H_0$.

Only $A_i(z)$ changes under this delay: it loses $x$, while all
other $A_j(z)$ remain unchanged.  Correspondingly, $m_i(a,z)$
decreases by one and all other $m_j(a,z)$ remain unchanged.
The inclusion in test $i$ stays true because $x\notin R_i$.
The left side of that test decreases by
$\min\{\lambda_H(x),\varepsilon_i\}\le\varepsilon_i$, whereas its
right side decreases by exactly $\varepsilon_i$.
Thus the same $H$ and cutoffs witness feasibility of $z'$.
But $L(a,z')=L(a,z)-1$, contradicting the defining minimum.
This proves \eqref{sel:changed-vertices} for every $a$ in the group.
Figure~\ref{fig:one-level-delay} illustrates why moving one vertex back
by one level changes only one of the feasibility tests.

\begin{figure}[ht]
\centering
\begin{tikzpicture}[x=1.2cm,y=1cm,>=Latex,font=\small]
\draw[->,gray] (0,0) -- (7,0) node[right]{color};
\foreach \x/\lab in {1/{1},3/{i},4/{i+1},6/{h}} {
  \draw (\x,-.08) -- (\x,.08);
  \node[below] at (\x,-.1) {$\lab$};
}
\fill[blue!55!black] (3,0) circle (2pt);
\fill[blue!55!black] (6,0) circle (2pt);
\draw[->,thick,blue!55!black] (5.95,.15) to[bend right=30]
 node[above] {chosen move} (3.05,.15);
\draw[->,thick] (3,-.7) -- node[below]{delay one level} (4,-.7);
\draw[dashed,gray] (3.5,-1.2) -- (3.5,.55);
\node[anchor=east] at (3.3,-1.5) {inside $A_i$};
\node[anchor=west] at (3.7,-1.5) {outside $A_i$};
\end{tikzpicture}

\caption{The one-level delay in the reconstruction argument. A vertex
originally colored $h>i$ has been moved to $i$; delaying it to $i+1$
removes it from $A_i$ alone. At the same time $m_i$ decreases by one,
lowering the right side of \eqref{sel:feasibility} by $\varepsilon_i$.
Its membership in every other cumulative set $A_j$ stays fixed.
The drawing is schematic; $h$ may equal $i+1$.}
\label{fig:one-level-delay}
\end{figure}

\medskip
\noindent\textit{Counting the original colorings.}
Given $z$ and the profile, an original coloring $a$ is determined by
its changed vertices of each color $i$ in $z$, together with their
original colors $h>i$.  By \eqref{sel:changed-vertices}, the number
of possibilities is at most
\begin{equation}
\label{sel:encoding-bound}
  \prod_{i=1}^s
    \binom{|R_i|}{t_i}
    \binom{t_i}{(n_{ih})_{h>i}},
\end{equation}
where the second factor is a multinomial coefficient.
The sets $R_i$ may overlap and may contain vertices of other colors
in $z$.  Counting arbitrary subsets of them only adds possibilities
to \eqref{sel:encoding-bound}; every actual coloring has a unique
encoding of the indicated form.

For $t>0$, \eqref{sel:large-weight-bound} and
$\binom{n}{k}\le(en/k)^k$ give
\begin{equation}
\label{sel:entropy-bound}
  \prod_{i=1}^s\binom{|R_i|}{t_i}
  \le
  \prod_{i:t_i>0}\left(\frac{e\,2^i t}{t_i}\right)^{t_i}
  \le \prod_{i=1}^s(e\,4^i)^{t_i}.
\end{equation}
To see the last inequality, apply weighted AM--GM with weights
$\alpha_i=t_i/t$ on the positive rows:
\[
  \prod_{i:t_i>0}
     \left(\frac{2^{-i}}{\alpha_i}\right)^{\alpha_i}
  \le \sum_{i:t_i>0}2^{-i}\le1.
\]
Raising this inequality to the power $t$ yields exactly the second
inequality in \eqref{sel:entropy-bound}.  Rows with $t_i=0$
contribute a factor of one.

For each pair in the group, its probability and cost satisfy the
exact identity
\begin{equation}
\label{sel:likelihood}
  p^tP(a)=P(z)
      \prod_{i=1}^s\left(\frac p{\pi_i}\right)^{t_i}
      \prod_{i<h}\pi_h^{n_{ih}}.
\end{equation}
All color probabilities are positive.  Combining
\eqref{sel:encoding-bound}--\eqref{sel:likelihood}, the contribution
of a fixed $z$ and profile to the right side of
\eqref{sel:cover-cost} is at most
\begin{equation}
\label{sel:profile-contribution}
  P(z)\prod_{i=1}^s
  \left[
     \left(\frac{e\,4^i p}{\pi_i}\right)^{t_i}
     \binom{t_i}{(n_{ih})_{h>i}}
     \prod_{h>i}\pi_h^{n_{ih}}
  \right].
\end{equation}
The witness and the sets $R_i$ can depend on the full profile, but
none of that dependence remains in this bound.

For fixed row sums $t_1,\ldots,t_s$, we can consequently sum
\eqref{sel:profile-contribution} over all nonnegative integer
profiles with those row sums, including profiles that do not
occur.  The multinomial theorem gives, for each row,
\[
  \sum_{\substack{(n_{ih})_{h>i}\ge0\\
                   \sum_{h>i}n_{ih}=t_i}}
      \binom{t_i}{(n_{ih})_{h>i}}
      \prod_{h>i}\pi_h^{n_{ih}}
  =\left(\sum_{h>i}\pi_h\right)^{t_i}\le1.
\]
Finally, summing over all colorings $z$ uses $\sum_zP(z)=1$.
Thus the right side of \eqref{sel:cover-cost} is at most
\begin{equation}
\label{sel:tuple-sum}
  \sum_{\substack{t_1,\ldots,t_s\ge0\\
                   t_1+\cdots+t_s>0}}
       \prod_{i=1}^s
         \left(\frac{e\,4^i p}{\pi_i}\right)^{t_i}
  =
  \sum_{\substack{t_1,\ldots,t_s\ge0\\
                   t_1+\cdots+t_s>0}}
       \prod_{i=1}^s\left(\frac e{64^i}\right)^{t_i},
\end{equation}
where both sums run over integer tuples.
Put $\rho=\sum_{i=1}^s e/64^i$.  For a fixed total
$\ell=t_1+\cdots+t_s$, each monomial in
\eqref{sel:tuple-sum} occurs in the expansion of $\rho^\ell$ with
a multinomial coefficient at least one.  Since
\[
  \rho<\frac e{63}<\frac1{21},
\]
we obtain
\[
  \text{right side of \eqref{sel:cover-cost}}
  \le\sum_{\ell\ge1}\rho^\ell
  =\frac{\rho}{1-\rho}<\frac1{20}.
\]
Together with \eqref{sel:cover-cost}, this proves that the bad
probability is at most $1/10$.

For every nonbad coloring the same $H$ satisfies all inequalities
in \eqref{sel:simultaneous-capture}.  Since the colors are integers,
\[
  \sum_{x\in X}a(x)\lambda_H(x)
  =1+\sum_{i=1}^s\bigl(1-\lambda_H(A_i(a))\bigr)
  \le1+\sum_{i=1}^s2^{-i}\le2.
\]
This proves \eqref{sel:average-color} and completes the proof.
\end{proof}
